\chapter[INTRODUÇÃO]{INTRODUÇÃO}

\section{Tema Delimitado}
O tema deste trabalho é a automação do controle de qualidade de rochas ornamentais por meio de um sistema hierárquico de Inteligência Artificial, denominado ARIA (Análise e Reconhecimento Inteligente de Anomalias). O ARIA é o \textit{pipeline} de visão computacional proposto e validado nesta pesquisa, concebido para ser o núcleo inteligente da plataforma industrial Hartheus, voltada ao controle de qualidade de chapas. Mais especificamente, a pesquisa aborda o desenvolvimento e a validação deste pipeline computacional que integra modelos de \textit{Deep Learning} em duas etapas: um classificador taxonômico baseado na arquitetura Xception (DeepStoneAI), responsável por identificar o domínio litológico da chapa; e um segmentador especialista, estruturado em uma arquitetura \textit{Teacher-Student}. Este segmentador utiliza o \textit{Segment Anything Model} (SAM) para geração de anotações de alta precisão (\textit{pseudo-labeling}) que, por sua vez, alimentam o treinamento de um modelo YOLO11 focado na detecção rápida de anomalias superficiais (fissuras, oxidações e manchas). A pesquisa tem aplicação direta no setor industrial do Espírito Santo — principal polo produtor e exportador do Brasil —, alinhando-se aos conceitos da Indústria 4.0.

\section{Problema de Pesquisa}
A inspeção de defeitos superficiais em chapas de rochas ornamentais é realizada de forma manual, dependendo da acuidade visual humana, o que a torna subjetiva, propensa a falhas e de difícil padronização. A alta variabilidade geológica das rochas dificulta o treinamento de modelos de IA generalistas, tornando a distinção entre características naturais (como veios minerais) e defeitos reais um desafio técnico de difícil parametrização. Além disso, a anotação manual de milhares de defeitos sutis para treinamento é um processo oneroso e limitante. Diante disso, formula-se a seguinte pergunta: a implementação do sistema ARIA, um pipeline hierárquico que utiliza classificação litológica prévia para calibrar uma arquitetura \textit{Teacher-Student} (SAM e YOLO11), é capaz de automatizar a geração de anotações (\textit{pseudo-labeling}) e superar a precisão de modelos generalistas na detecção de anomalias em rochas ornamentais?

\section{Justificativa}
Economicamente, o projeto justifica-se pela importância de padronizar o controle de qualidade em um setor estratégico para as exportações brasileiras, garantindo uma precificação mais justa e reduzindo perdas industriais. Do ponto de vista tecnológico e acadêmico, o trabalho propõe uma solução inovadora para o gargalo da escassez de dados anotados na indústria, preenchendo uma lacuna identificada na literatura sobre a integração hierárquica entre modelos de classificação e segmentação. Ao introduzir a injeção de conhecimento de domínio na calibração do modelo SAM (Professor), o sistema automatiza a criação de \textit{datasets} de alta qualidade para o treinamento do modelo YOLO11 (Aluno). Isso permite criar um ciclo contínuo de inferência rápida em tempo real, alinhando a manufatura do setor de rochas ornamentais aos paradigmas da Indústria 4.0 de forma viável e escalável.

\section{Objetivo Geral}
Desenvolver, otimizar e validar o sistema ARIA, um pipeline hierárquico de Inteligência Artificial para o controle de qualidade de rochas ornamentais, integrando a classificação taxonômica (DeepStoneAI) a um segmentador especialista em arquitetura \textit{Teacher-Student}, visando automatizar a detecção de anomalias superficiais e reduzir o índice de falsos positivos em comparação a abordagens generalistas.

\section{Objetivos Específicos}
\begin{itemize}
    \item Constituir e pré-processar um banco de imagens industriais de rochas ornamentais com variabilidade de classes e anomalias sutis.
    \item Implementar o módulo de classificação global (DeepStoneAI) capaz de identificar e rotear a rocha para o fluxo de análise correspondente.
    \item Desenvolver e calibrar o modelo SAM (Professor) utilizando injeção de conhecimento de domínio (arquivos de \textit{prompts} específicos por litologia) para a geração automatizada de \textit{pseudo-labels} de anomalias.
    \item Treinar o modelo YOLO11 (Aluno) utilizando as anotações geradas pelo SAM, criando um motor de inferência rápida e eficiente para o ambiente de produção.
    \item Avaliar o desempenho do sistema utilizando métricas quantitativas (como mAP e IoU) em conjunto com análises qualitativas, validando a capacidade do modelo de superar a subjetividade da anotação humana inicial.
\end{itemize}

\chapter[FUNDAMENTAÇÃO TEÓRICA]{FUNDAMENTAÇÃO TEÓRICA}

Para sustentar a proposta do sistema ARIA, este referencial teórico aborda as bases organizacionais da adoção tecnológica e os conceitos de \textit{Deep Learning} que fundamentam a classificação hierárquica e a segmentação visual.

\section{Automação e Sistemas de Informação no Setor de Rochas}
A inspeção de qualidade no setor de rochas ornamentais é uma etapa crítica, tradicionalmente dependente de operadores humanos. Por ser um processo visual, ele é inerentemente subjetivo, suscetível à fadiga do trabalhador e propenso a inconsistências. Para mitigar esses problemas, a adoção de sistemas baseados em Inteligência Artificial atua como uma solução estratégica para modernizar o fluxo de trabalho. A implementação em ambientes industriais envolve dinâmicas organizacionais explicadas por Teorias de Sistemas de Informação:
\begin{itemize}
    \item \textbf{Modelo de Aceitação de Tecnologia (TAM):} Indica que a eficácia da introdução de uma inovação técnica depende de como os operadores percebem a utilidade do sistema (redução da carga de trabalho repetitivo) e sua facilidade de uso.
    \item \textbf{Teoria da Difusão de Inovações:} Explica que a aceitação requer a demonstração clara de vantagens relativas, como o aumento da eficiência e da padronização.
    \item \textbf{Teoria Sociotécnica:} Ressalta que o sucesso provém da harmonização entre os aspectos técnicos (a Inteligência Artificial) e sociais (a qualificação humana e organização do ambiente).
\end{itemize}

\section{Redes Neurais Convolucionais e Classificação Taxonômica}
Para viabilizar o controle automatizado, a Visão Computacional baseada em Redes Neurais Convolucionais (CNNs) atua na extração de informações de imagens por meio de filtros que evidenciam texturas e formas. Em um pipeline hierárquico, a identificação do domínio de conhecimento (a variedade litológica da rocha) é a etapa primordial. Para esta tarefa, a arquitetura Xception apresenta alta eficiência através de convoluções separáveis em profundidade (\textit{depthwise separable convolutions}), reduzindo o custo computacional e garantindo alta acurácia na etapa de roteamento do sistema (DeepStoneAI).

\section{Modelos de Fundação e o Segment Anything Model (SAM)}
Nos últimos anos, a Inteligência Artificial avançou para o paradigma dos Modelos de Fundação (\textit{Foundation Models}) — redes neurais massivas treinadas em conjuntos de dados em larga escala, capazes de realizar tarefas genéricas com alta precisão (capacidade \textit{zero-shot}). Neste contexto, destaca-se o \textit{Segment Anything Model} (SAM). Diferente de arquiteturas tradicionais, o SAM é fundamentado na engenharia de \textit{prompts}, permitindo que o modelo seja guiado por coordenadas, caixas delimitadoras ou textos para isolar objetos específicos em uma imagem. No domínio de rochas ornamentais, o SAM permite a injeção de conhecimento de domínio: ao calibrar os \textit{prompts} com os limiares visuais específicos de cada litologia, evita-se que o modelo confunda características geológicas naturais (veios) com defeitos comerciais.

\section{Arquitetura Teacher-Student e Pseudo-labeling}
O principal gargalo para a implementação de Visão Computacional na indústria é a escassez de dados anotados manualmente (\textit{Ground Truth}). Para contornar essa limitação, a literatura adota a transferência de conhecimento por meio da arquitetura \textit{Teacher-Student} acoplada à técnica de \textit{pseudo-labeling}. Nesta abordagem, um modelo robusto e de alto custo computacional atua como Professor (neste trabalho, o SAM), processando imagens brutas de forma \textit{offline} para gerar máscaras de segmentação automatizadas. Essas anotações simuladas, denominadas \textit{pseudo-labels}, são então utilizadas para supervisionar e treinar um modelo Aluno, que é estatisticamente menor, porém otimizado para inferências rápidas.

\section{Inferência em Tempo Real com YOLO}
No ambiente industrial de manufatura contínua, onde prevalecem hardwares com restrição de memória (\textit{Edge Computing}), a velocidade de inferência é prioritária. O modelo Aluno escolhido é a arquitetura YOLO (\textit{You Only Look Once} - versão 11), que se destaca por processar extração de características (Backbone) e predição de máscaras (Head) em uma única passagem pela rede, atingindo altas taxas de quadros por segundo (FPS). 
Essa eficiência viabiliza a inferência em tempo real na linha de produção, requisito central para a aplicação industrial proposta.

\chapter[METODOLOGIA]{METODOLOGIA}

\section{Metodologia}

\subsection{Tipo de Pesquisa}
O presente trabalho classifica-se como uma pesquisa aplicada, de natureza experimental e abordagem mista (quantitativa e qualitativa). É aplicada pois busca resolver um problema prático do setor industrial capixaba (a subjetividade e o custo da inspeção manual de chapas de rochas ornamentais). Caracteriza-se como experimental pois envolve a manipulação de variáveis, arquiteturas de redes neurais e limiares de hiperparâmetros para observar seus efeitos nas métricas de acurácia, precisão e velocidade de inferência.

\subsection{Procedimentos Metodológicos}
O desenvolvimento do sistema ARIA será conduzido em etapas sequenciais, garantindo a validação de cada módulo antes da integração final:

\begin{enumerate}
    \item \textbf{Coleta e Pré-processamento de Dados:} Constituição de um banco de imagens industriais reais e customizadas, garantindo a representatividade de diferentes litologias e o desbalanceamento natural de classes (anomalias sutis vs. chapas perfeitas).
    \item \textbf{Implementação do Módulo de Classificação (DeepStoneAI):} Treinamento e validação da rede Xception para classificar o domínio litológico da chapa, estabelecendo a base para a restrição de domínio da próxima etapa.
    \item \textbf{Calibração do Professor e Geração de \textit{Pseudo-labels}:} Injeção de conhecimento de domínio no modelo \textit{Segment Anything Model} (SAM). Será utilizado um arquivo de configuração estruturado (\texttt{rock\_prompts.json}) para definir parâmetros e restrições visuais específicas de cada rocha. O SAM processará o \textit{dataset} de forma \textit{offline}, gerando as anotações automatizadas. Nesta versão, todas as anomalias são tratadas como uma única classe genérica (rotulagem binária), uma vez que a distinção semântica entre tipos de anomalia depende da qualidade dos \textit{embeddings} de texto, ainda não validada de forma independente; a rotulagem multiclasse fica registrada como extensão futura.
    \item \textbf{Treinamento do Aluno (YOLO11):} Alimentação do motor de inferência rápida utilizando os \textit{pseudo-labels} gerados na etapa anterior. O modelo YOLO será treinado para maximizar a taxa de quadros por segundo (FPS) e a assertividade na detecção.
    \item \textbf{Validação:} Avaliação quantitativa utilizando métricas como \textit{Mean Average Precision} (mAP) e \textit{Intersection over Union} (IoU), complementada pela validação qualitativa das anotações por especialistas do setor de rochas ornamentais.
\end{enumerate}

\subsection{Protocolo Experimental}
O experimento central avalia a hipótese da especialização hierárquica por meio da comparação entre dois \textit{pipelines} treinados sobre a mesma arquitetura YOLO11-seg. No \textit{pipeline} especialista (ARIA), o classificador Xception identifica o tipo litológico e seleciona os \textit{prompts} calibrados correspondentes, resultando em \textbf{45 modelos especializados} --- um por tipo de rocha, treinado exclusivamente nas anotações daquela litologia. No \textit{pipeline} \textit{baseline} (generalista), o estágio de classificação é ausente: o SAM utiliza uma configuração genérica de \textit{prompts} para todas as rochas e um \textbf{único modelo} YOLO é treinado com as anotações de todos os tipos. A comparação isola a variável da hierarquia --- a quantidade de modelos (45 \textit{vs.} 1) e a especificidade das anotações ---, permitindo medir o ganho da abordagem especialista nas métricas de mAP e IoU.

\section{Ferramentas e Tecnologias Utilizadas}
Para a condução dos experimentos e desenvolvimento do pipeline computacional, será adotada a seguinte \textit{stack} tecnológica:
\begin{itemize}
    \item \textbf{Linguagem e Bibliotecas:} Desenvolvimento primário em Python, utilizando o \textit{framework} PyTorch para a construção, treinamento e inferência dos tensores e modelos de \textit{Deep Learning}.
    \item \textbf{Aceleração de Hardware:} Utilização da plataforma CUDA (NVIDIA) para viabilizar o processamento paralelo em GPU, reduzindo drasticamente o tempo de treinamento dos modelos e permitindo a inferência do SAM.
    \item \textbf{Validação Qualitativa:} Amostras dos \textit{pseudo-labels} gerados serão submetidas à avaliação de especialistas do setor de rochas ornamentais, verificando se as segmentações identificam anomalias reais do ponto de vista industrial.
\end{itemize}

\section{Elemento Visual: Arquitetura do Sistema}
A Figura \ref{fig:arquitetura_hartheus} ilustra o fluxo de processamento e a hierarquia do sistema ARIA, detalhando a transição desde a classificação taxonômica até a segmentação de anomalias pelas arquiteturas \textit{Teacher-Student}.

% ATENÇÃO: Substitua 'nome_da_sua_imagem.png' pelo nome do arquivo que você fará upload no Overleaf.
\begin{figure}[htpb]
    \centering
    \caption{Pipeline hierárquico do sistema ARIA: Integração do DeepStoneAI e segmentador especialista (SAM e YOLO11).}
    \label{fig:arquitetura_hartheus}
\end{figure}

\section{Cronograma de Execução}
O planejamento das atividades para a concretização da pesquisa está distribuído ao longo dos meses de desenvolvimento, conforme a Tabela \ref{tab:cronograma}.

\begin{table}[htpb]
    \centering
    \caption{Cronograma de execução das atividades do projeto.}
    \vspace{0.2cm}
    % O segredo está aqui: trocamos 'l' por 'p{6cm}'
    \begin{tabular}{|p{6cm}|c|c|c|c|c|c|}
        \hline
        \textbf{Atividades} & \textbf{Mês 1} & \textbf{Mês 2} & \textbf{Mês 3} & \textbf{Mês 4} & \textbf{Mês 5} & \textbf{Mês 6} \\ \hline
        Levantamento Bibliográfico e Estruturação & X & X & & & & \\ \hline
        Coleta e Pré-processamento do Dataset & & X & X & & & \\ \hline
        Implementação e Teste do DeepStoneAI & & & X & X & & \\ \hline
        Calibração do SAM e Geração de Labels & & & & X & X & \\ \hline
        Treinamento e Avaliação do YOLO11 & & & & & X & X \\ \hline
        Escrita do Documento e Defesa & X & X & X & X & X & X \\ \hline
    \end{tabular}
    \label{tab:cronograma}
\end{table}

\chapter[RESULTADOS INICIAIS]{RESULTADOS INICIAIS}

Este capítulo apresenta as primeiras evidências práticas do desenvolvimento do sistema, em consonância com a proposta descrita na metodologia. O foco atual encontra-se na etapa do segmentador especialista, mais especificamente na atuação do modelo \textit{Segment Anything Model} (SAM) como Professor, responsável pela geração automatizada de anotações (\textit{pseudo-labeling}). Os resultados a seguir foram obtidos na rocha \textit{ice leke}, primeiro litótipo a ter seus \textit{prompts} e limiares de confiança calibrados e validados no arquivo de configuração \texttt{rock\_prompts.json}.

\section{Geração de Anotações pelo Professor (SAM)}
A partir de uma imagem industrial real da chapa, o SAM foi guiado por \textit{prompts} de linguagem natural calibrados para a litologia (injeção de conhecimento de domínio), segmentando as anomalias superficiais sem qualquer treinamento prévio específico para rochas ornamentais (capacidade \textit{zero-shot}). A Figura \ref{fig:resultado_combined} contrapõe a imagem original às máscaras de segmentação sobrepostas, geradas de forma totalmente automatizada.

\begin{figure}[htpb]
    \centering
    \caption{Segmentação automatizada de anomalias gerada pelo SAM na rocha \textit{ice leke}: imagem original (a) e máscaras de anomalias sobrepostas (b).}
    \label{fig:resultado_combined}
    \subfloat[Imagem original da chapa\label{subfig:original}]{\includegraphics[width=0.45\textwidth]{sample_ice_leke_original.jpg}}
    \hfill
    \subfloat[Anotações geradas pelo SAM\label{subfig:combined}]{\includegraphics[width=0.45\textwidth]{sample_ice_leke_combined.jpg}}
\end{figure}

\section{Decomposição por Tipo de Anomalia}
Cada \textit{prompt} calibrado atua sobre um tipo específico de irregularidade. A Figura \ref{fig:resultado_decomposto} apresenta as máscaras isoladas por categoria — fissuras (\textit{crack}), veios minerais (\textit{vein}), manchas (\textit{Stain}) e áreas escuras (\textit{Dark patches}) —, evidenciando a capacidade do modelo de distinguir diferentes feições visuais a partir de descrições textuais.

\begin{figure}[htpb]
    \centering
    \caption{Decomposição das anotações do SAM por tipo de anomalia na rocha \textit{ice leke}.}
    \label{fig:resultado_decomposto}
    \subfloat[\textit{crack} (fissuras)\label{subfig:crack}]{\includegraphics[width=0.45\textwidth]{sample_ice_leke_crack.jpg}}
    \hfill
    \subfloat[\textit{vein} (veios)\label{subfig:vein}]{\includegraphics[width=0.45\textwidth]{sample_ice_leke_vein.jpg}}
    \\[0.4cm]
    \subfloat[\textit{Stain} (manchas)\label{subfig:stain}]{\includegraphics[width=0.45\textwidth]{sample_ice_leke_stain.jpg}}
    \hfill
    \subfloat[\textit{Dark patches} (áreas escuras)\label{subfig:darkpatches}]{\includegraphics[width=0.45\textwidth]{sample_ice_leke_darkpatches.jpg}}
\end{figure}

\section{Análise Preliminar}
Os resultados confirmam a viabilidade da etapa central da proposta: o modelo de fundação, sem treinamento dedicado ao domínio, é capaz de segmentar anomalias relevantes quando guiado por \textit{prompts} calibrados por litologia. As máscaras resultantes são exportadas como polígonos no formato de segmentação do YOLO (coordenadas normalizadas), constituindo o conjunto de \textit{pseudo-labels} que alimentará o treinamento do modelo Aluno (YOLO11). Trata-se de uma validação \textbf{qualitativa e preliminar}: das 45 litologias do conjunto de dados, a \textit{ice leke} é a primeira plenamente calibrada, e a calibração das demais segue em andamento. As métricas quantitativas (mAP e IoU) serão obtidas após o treinamento do Aluno, em etapa posterior do cronograma.